\subsection{Language Models}
Before describing the literature of \ac{ie} and related tasks, we first introduce what \acl{ai} is and the recent advances in \ac{nlp}, with special focus on \acp{llm} and on their relation with \acp{kb}.

The question ``Can a machine think?'' dates back to the 1950 when the Turing test was proposed~\cite{turing1950computing}. A computer passes the test if, after asking a series of written questions, a human evaluator cannot distinguish whether the answers are generated by a person or by the computer~\cite{}.
\Acf{ai} refers to the abilities a machine may require to emulate human behavior, including \acl{nlp} to understand and generate human language, \acl{kr} to organize and store knowledge, automated reasoning to answer questions and infer new conclusions, \acl{ml} to adapt and detect patterns, as well as computer vision and speech recognition to perceive the environment, and robotics to manipulate objects and navigate the world~\cite{russell2020artificial}.

A precise definition of an AI system has been established by \textcite{ai_act}---a European regulation designed to ensure \acl{ai} systems are trustworthy and human-centric, promoting innovation while protecting health, safety, fundamental rights, and the environment~\cite[Article 1]{ai_act}.

\begin{definition}[Artificial Intelligence]
``\emph{\acs{ai} system} means a machine-based system that is designed to operate with varying levels of autonomy and that may exhibit adaptiveness after deployment, and that, for explicit or implicit objectives, infers, from the input it receives, how to generate outputs such as predictions, content, recommendations, or decisions that can influence physical or virtual environments''~\cite[Article 3]{ai_act}.
\end{definition}

\todo{collegare meglio, e.g. in this thesis I worked more on text,...}

\emph{\Acf{nlp}} refers to the processing and understanding of a sequence of words in natural language by a computer~\cite{bates1995models}.
Its research dates back to the 1950s, initially focusing on foundational tasks such as machine translation and question answering~\cite{bates1995models}.\todo{definition of QA here?}
Over the years, \ac{nlp} has evolved significantly, with early approaches relying on rule-based systems and statistical methods, while more recent advancements have been driven by machine learning techniques, and in particular by deep learning ones~\cite{entityorientedsearch2018}.

A \emph{deep learning neural network} is a network composed of multiple layers (hence \emph{deep}) of small computing units, called \emph{neurons}, where each unit processes an input vector and returns a single output value~\cite{jm3}.
%
\Acf{dl} is typically data-driven~\cite{jm3}, i.e., \ac{dl} neural networks are trained on a training dataset, through the minimization of a loss function, representing the error of the network predictions on the training set~\cite{hastie2009elements}.

Recently, the introduction of transformer architectures~\cite{vaswani_transformers_2017} has revolutionized the field, establishing them as the dominant approach for building large language models, i.e., models pretrained on massive data that can be used or adapted for downstream tasks~\cite{gu2024mamba}.
\begin{definition}
    [Language Model] A \acf{lm} is a probability distribution over a vocabulary of tokens, conditioned on an input sequence, and used to predict the next token~\cite{lm_java_methods_2023,jm3}.
\end{definition}
\todo{and a large language model?}
% formalization
We speak about \emph{tokens} instead of words because modern \ac{lm} do not operate at the word level but rather on subword units, which can be whole words, word parts, symbols, or characters~\cite{comprensive_overview_llm_2025}. The use of a token-based vocabulary, instead of a word-based or character-based one, derives from the need to balance vocabulary size and coverage. Word-based vocabularies suffer from the issue of out-of-vocabulary words. Character-based vocabularies, on the other hand, can lead to longer sequences and increased computational costs~\cite{sun-etal-2023-characters}. Token-based tokenizers, such as Byte-Pair Encoding (BPE)~\cite{sennrich-etal-2016-neural}, can represent frequent words using a single token, still allowing the representation of rare or novel words by combining multiple tokens.\todo{verify correctness}

This definition of \ac{lm} is in-line with \emph{casual language modeling (CLM)}, that predicts the next token given the previous ones~\cite{jain-etal-2023-cl-for-causal-lm,jm3}, and these models are also reffered to as \emph{autoregressive models}~\cite{jm3}, i.e. ``models regressing the outcomes on previous values of the same time series''~\cite{regis2022_autoregressive}, where in \ac{nlp} we have a sequence of tokens instead of a time series.
\todo{causal lm objective and dl training and formulation}

\subsubsection{Causal Language Modeling}

Given a sequence of tokens $x = [x_0, x_1, \ldots, x_|x|]$, a causal \ac{lm} estimates the probability of the token $x_i$ as
\begin{equation}
    P(x_i \mid x_1, x_2, \ldots, x_{i-1}) = P(x_i \mid x_{<i}),
\end{equation}
while the probability of the whole sequence is commonly estimated by the factorization~\cite{jain-etal-2023-cl-for-causal-lm,jm3}
\begin{equation}
    P(x) = \prod_{i=1}^{|x|} P(x_i \mid x_{<i}).
\end{equation}
Deep-learning based causal \ac{lm} are trained via maximum likelihood estimation~\cite{jain-etal-2023-cl-for-causal-lm}, i.e., by maximizing the log-likelihood of the training data~\cite{MLE_tutorial}, that corresponds to $N$ sequences of tokens $[x^0, x^1, \ldots, x^|N|]$. The training loss is thus
\begin{equation}
    \mathcal{L}_{CLM} = - \frac{1}{N} \sum_{j=1}^{N} \log P(x^j) = - \sum_{j=1}^{N} \sum_{i=1}^{|x^j|} \log P(x_i^j \mid x_{<i}^j).
\end{equation}

\Aclp{clm}, that estimates the probability distribution for the next token given the previous ones, are used for text generation by iteratively sampling the next token from the distribution (e.g., greedily selecting the most probable token) and appending it to the input sequence~\cite{jm3}. 

\subsubsection{Masked Language Modeling}

Another paradigm is \emph{masked language modeling}, where some tokens in the input sequence are randomly \emph{masked}, i.e., replaced with a special token, and the model is trained to predict the original tokens based on the bidirectional (left and right) surrounding context~\cite{devlin-etal-2019-bert}.

Formally, given a sequence of tokens $x = [x_1, x_2, \ldots, x_{|x|}]$, a random subset of positions 
$M \subseteq \{1, \ldots, |x|\}$ is selected for masking, yielding a corrupted input sequence $\tilde{x}$. 
For each masked position $i \in M$, a masked \ac{lm} estimates the probability of the original token $x_i$ conditioned on its left and right context:
\begin{equation}
    P(x_i \mid \tilde{x}_{\setminus i}) = P(x_i \mid x_{<i}, x_{>i}),
\end{equation}
where $x_{<i}$ and $x_{>i}$ denote the tokens preceding and following position $i$, respectively. 
The training loss, corresponding to maximum likelihood estimation over the masked positions of $N$ sequences, is
\begin{equation}
    \mathcal{L}_{MLM} = - \frac{1}{N} \sum_{j=1}^{N} \sum_{i \in M^j} \log P(x_i^j \mid x_{<i}^j, x_{>i}^j).
\end{equation}

While \acp{clm} are suited for text generation, \acp{mlm} are typically used for sequence labeling or token classification tasks, taking advantage of the context from both sides~\cite{jm3}. Examples of such tasks include \acf{pos} tagging, where each token is assigned a part-of-speech category such as \texttt{Noun}, \texttt{Verb}, or \texttt{Adj}, and \acf{ner}---described in \todo{sect ner}---where each token is assigned a label indicating whether it is part of a named entity and its type, such as \texttt{Person}, \texttt{Location}, or \texttt{Organization}~\cite{jm3}.

\subsubsection{Transformers}
\emph{Transformers} have been introduced in 2017 and evaluated on the machine translation, a language modeling task, demonstrating superior performance and higher parallization at training-time compared to previous architectures~\cite{vaswani_transformers_2017}.
\
todo{say why machine translation is language modeling}
The transformer uses an encoder-decoder architecture and it is cabable of modeling sequences of tokens relying on the \emph{attention mechanism}~\cite{vaswani_transformers_2017}.\todo{describe encoder decoder, attention and self attention}.
\todo{briefly introduce representation learning, word embeddings and maybe knowledge grah embeddings. jm3 or repr learn 4 nlp}
An \emph{encoder} processes tokens as input and transforms each into a vector representation, commonly referred to as an embedding~\cite{jm3}.
% To place this concept in context, it can be connected to the broader field of \emph{representation learning}, which---with respect to text---aims at automatically learn useful representations of words, sentences, or documents~\cite{jm3}.

\paragraph{Vector Semantics and Representation Learning}
To place the concept of \emph{embedding} in context, it can be connected to the fields of \emph{representation learning} and \emph{vector semantics}.
\emph{Representation learning} provides methods for learning useful representations from data, including text (words, sentences, documents)~\cite{bengio_representation_learning_2013}.\todo{ragiona se togliere repr learning}
With respect to \emph{vector semantics}, its origin can be traced to the 1950s with the \emph{distributional hypothesis}~\cite{Harris01081954}---i.e., words occurring in similar contexts tend to have similar meanings~\cite{jm3}---and the idea of using points in a three-dimentional space for representing words connotation~\cite{osgood1957}. Words like ``car'' and ``automobile'' are synonyms, but their \say{similarity} cannot be captured via string similarity (e.g., edit distance~\cite{jm3}). Also ``cats'' and ``dogs'', or ``coffee'' and ``cup'' can be considered similar or related.\todo{paraphrase the examples}
The intuition of vector semantics is to assign each word an \emph{embedding}, i.e., a vector for representing a word in a multi-dimentional semantic space, where similarities between embeddings would reflect the similarities of words' meanings~\cite{jm3}.

Different techniques have been proposed for obtaining word embeddings. In 2013, \textcite{mikolov2013} introduced \emph{word2vec}, which, in the case of the CBOW (Continuous Bag of Words) architecture, learns dense semantic embeddings by training a model to predict a target word from its surrounding context~\cite{mikolov2013}. \emph{Dense} embeddings refer to low-dimentional vectors, in constrast with very long \emph{sparse} vectors, such as the ones one-hot vectors, whose size corresponds to the length of the entire vocabulary (often between 10,000 and 50,000 words)~\cite{jm3}.
%
A standard illustration of the utility of word embeddings is provided by the following analogy:
\begin{equation}
\vec{\textsc{King}} - \vec{\textsc{Man}} + \vec{\textsc{Woman}} \approx \vec{\textsc{Queen}}.
\end{equation}
Subtracting the embedding of ``man'' from ``king'', then adding ``woman'' produces a vector close to ``queen'', demonstrating some kind of  \say{vector-based reasoning}~\cite{mikolov-etal-2013-linguistic}.

However, assigning a single embedding to each word in the vocabulary is insufficient for polysemous words, such as ``mouse'' or ``bank'', which have multiple distinct senses. For instance, ``mouse'' can refer either to the small rodent or to the device used to control a computer, depending on the context~\cite{jm3}.
\todo{this is why people introduced contextual embeddings and these embedding are indeed often obtained with transformers. then go on with transformers encoders that use the context weighted by attention to enrich the embedding of the tokens layer by layer}

On the other hand, transformers are able to compute \emph{contextual embeddings}, i.e., dense vectors representing the meaning of a token in the context~\cite{jm3}.\todo{piture from jm3 fig 9.6}
%
\begin{figure}[h!]
    \centering
    \includegraphics[width=0.8\textwidth]{images/contextual_embeddings.png}
    \caption{Contextual embeddings of the word ``die'' calculated with BERT~\cite{devlin-etal-2019-bert}, projected in two dimensions. Figure from~\textcite{cohen_2019_contextual_embeddings}.}
    \label{fig:contextual_embeddings}
\end{figure}
%
\Cref{fig:contextual_embeddings} shows the contextual embeddings computed with BERT~\cite{devlin-etal-2019-bert}, an encoder-only transformer, for the word ``die'', used in different contexts with different meanings, both in English and in German. The embeddings are positioned in semantic clusters according to the meaning of ``die'': \textsc{die (German)}, \textsc{die (death)}, and \textsc{die (dice)}.


\paragraph{The Transformer Architecture}

% \begin{figure}[h!]
%     \centering
%     \includegraphics[height=10cm]{images/Transformer_mt.drawio.pdf}
%     \caption{The transformer encoder-decoder stack (applied to machine translation, English to Italian).}
%     \label{fig:transformer_encoder_decoder}
% \end{figure}
%
Turning back to the transformer architecture, it is composed of a stack of encoders and a stack of decoders, as depicted in \Cref{fig:transformer_all}.
%
Given an input sentence of $N$ tokens, the sentence is divided in tokens, or \emph{tokenized}, then each token is assigned a static embedding from an embedding matrix $E \in \mathbb{R}^{|V| \times d}$ where $V$ is the vocabulary of the model and $d$ the depth or number of dimensions of the embeddings.
Subsequently, the token embeddings are processed by encoder and decoder blocks that contextualize them, by incorporating the meaning of contextual tokens. \emph{Attention} is the mechanism that weights and aggregates the token embeddings~\cite{jm3}.
%
\begin{figure}[ht!]
    \centering
    \includegraphics[width=.9\textwidth]{images/Transformer_all.drawio.pdf}
    \caption{Overview of the transformer architectures.}
    \label{fig:transformer_all}
\end{figure}
\todo{verify starting token}
includeonly

The original encoder-decoder transformer~\cite{vaswani_transformers_2017}, depicted in \Cref{fig:transformer_all}(a), targets machine translation. For example, given the English sentence ``Milan is in Italy'', the expect output in Italian is ``Milano è in Italia''. 
The input tokens in English are provided to the encoder (bottom left) and the decoder (bottom right) receives a special starting token ``<start>''. In the first forward pass the model generates ``Milano''. At the next step, the encoder receives the same input\todo{cached? and caching also for generation}, while the decoder receives as input the already generated sequence of tokens in Italian, ``<start>'' and ``Milano'', to generate ``è'' (the Italian verb to be). This iterative process continues until the model generates an end-of-sequence token or reaches maximum generation length.\todo{valuta se cambiare immagine o farne una e non citare jalammar. fai immagine (divisa in 3) encoder-decoder (machine translation), encoder-only (NER), decoder-only (text generation) with zoom on encoder and decoder block showing attention and ffn}

\todo{speak of encoder-only and decoder-only. who is better for what}

Both encoder and decoder receive as input an embedding for each token and produce another embedding as output. Formally, let $X \in \mathbb{R}^{N \times d}$ be the input matrix of shape $N \times d$. After processing $X$ the first encoder produces $X^i \in \mathbb{R}^{N \times d}$, a matrix of the same shape of $X$. This process is reapeated for every encoder block, progressively enriching the tokens' embeddings with contextual information weighted by the attention mechanism. 

As visible in \Cref{fig:transformer_all}(a), both the encoder and decoder modules include a \emph{self-attention} layer and the decoder also has a \emph{cross-attention} layer.\todo{valuta se muovere dopo e aggiungi l'idea che i layer man mano arricchiscono le rappresentazioni con il contesto. parla di cross attention e magari cita bengio}
Self-attention can be considered as a mechanism for constructing contextual representations of a token's meaning by focusing on and combining information from nearby tokens~\cite{jm3}. Indeed, tokens in a sequence attend to the same sequence.
The cross-attention, instead, is perfomed between different sequences. For example in machine translation, the decoder, while generating in the target language, selects which parts of the source sentence (in the source language) to pay attention to~\cite{bahdanau2015neural}. In the transformer~\cite{vaswani_transformers_2017}, each decoder layer attends to the final representations from the encoder stack.

For further details of the attention mechanism and its implementation, refer to \textcite{jm3}.\todo{scegli dove metterlo e fallo anche con gli altri libri (balog)}

The \emph{encoder-decoder} transformer for machine translation features a language modeling head, as visible in \Cref{fig:transformer_all}(a) that estimates the probability $P(x_i \mid y,x_{<i})$ where $x$ are the tokens in the target language, while $y$ in the source language. Using both encoder and decoder layers allow the model to use two different vocabularies, one for each language~\cite{jm3}. This model is usually trained in a supervised settings on tranlational datasets~\cite{}.\todo{vaswani but maybe something else}

Besides the encoder-decoder architecture for machine translation, two relevant modifications are the \emph{encoder-only} and the \emph{decoder-only} transformers. The \emph{encoder-only}, whose architecture is schetched in \Cref{fig:transformer_all}(b), is especially used for embedding or sequence classification tasks, such as document retrieval with dense vector and named entity recognition~\cite{jm3,devlin-etal-2019-bert}, and it is trained with \ac{mlm} for learning to create useful token representations.\todo{due parole su document retrieval. controlla architettura dei moderni modelli di embedding}

\todo{decoder only}
\emph{Decoder-only} models can be used used for language modeling\todo{ensure lm refers to clm}, as visible in \Cref{fig:transformer_all}(c), but they use the same vocabulary for input and output tokens, differently from encoder-decoder models that are used for machine translation or even speech recognition---where the input represents speech~\cite{jm3}.

Today, the most popular architecture for \acp{llm} is the decoder-only~\cite{zhang2025encoderdecodergemmaimprovingqualityefficiency,qorib-etal-2024-decoder}, although recent work demostrates interest in the potential of the encoder-decoder architecture~\cite{zhang2025encoderdecodergemmaimprovingqualityefficiency}.

A core component of the success of transformers and, is \emph{transfer learning}, i.e., ``the transfer of knowledge from a related task that has already been learned''~\cite{torrey2010transfer}, which is often achieved via pre-training on huge corpora with self-supervised objectives, such as \ac{mlm} or \ac{clm}, that do not require human labeling, to subsequently fine-tune the models parameters to downstream tasks~\cite{jm3,devlin-etal-2019-bert,radford2018improving,Raffel2020t5}.
\textcite{devlin-etal-2019-bert}, \textcite{radford2018improving}, indeed, demonstrate that rich pre-training on generic language tasks, such as \acl{mlm} and \ac{clm}, is useful also for downstream tasks, e.g., \ac{ner} and \ac{qa}. Furthermore, \textcite{Raffel2020t5} and \textcite{radford2019unsupervised_learners}, respectively, have achieved state-of-the-art performance on several language understanding tasks by mapping them in a text-to-text format, and have demonstrated the ability of \acp{lm} in a zero-shot setting. In this context, \emph{zero-shot} refers to using a model for a task for which it has not fine-tuned for, by providing a natural language description of the task~\cite{browngpt3}.

Later, by scaling the number of model parameters, \acfp{llm}, such as GPT-3~\cite{browngpt3} counting 175 billions of parameters, researchers have been able to achieve surprising results in zero-shot and few-shot---like zero-shot but also providing few natural language examples of the task---settings, in some cases even competitive with state-of-the-art fine-tuning approaches~\cite{browngpt3}.
These improvements enabled user to \emph{prompt} or instruct language models to do specific task, eventually providing one-shot or few-shot examples to better take advantage of their \emph{\acf{icl}} capability~\cite{jm3}. A \emph{prompt}, indeed, is a textual istruction given the user gives to a model for accomplishing a specific task. \emph{\Ac{icl}} refers to learning a task or improving the performance on that task without updating model parameters, e.g., with few-shot prompting~\cite{jm3}.

However, increasing models' size in terms of number of parameters have not been enough to improve their behaviour at following users' instruction~\cite{instruct_gpt_2022}. Models, for instance, may generate unhelpful, toxic, or hallucinated~\cite{maynez-etal-2020-faithfulness-hallucination,ji_etal-2023-survey-hallucination}---i.e., incoherent, nonsensical, or unfaithful to the provided input---content. For this reason, techniques to \emph{align} models to user's intent have been applied to \acp{llm}, including some based on supervised fine-tuning with cureated examples~\cite{zhou2023lima} and on reinforcement learning, such as \emph{\acf{rlhf}}~\cite{instruct_gpt_2022}.
Also methods for improving reasoning in \acp{llm} have been studied. \Acf{cot}~\cite{wei2022chainofthought} consists in decomposing a complicated reasoning task in multiple intermediate steps, which are solved to reach the final solution. More recently, so-called \emph{\acfp{lrm}}~\cite{tie2025surveyposttraininglargelanguage} have been trained to generate long \aclp{cot}~\cite{openai2024openaio1card,deepseekai2025deepseekr1incentivizingreasoningcapability}.

\todo{maybe add some definitions}



For a broader and more detailed perspective on the history of \acp{llm} refer to \textcite{Zhou2024_from_bert_to_chat}\todo{use nlp description from this paper} and \textcite{tie2025surveyposttraininglargelanguage}.


\todo{efficiency of models}


\todo{instruction-following. prompt-based. generalization to lot of tasks without fine-tuning. transfer learning. in-context learning}

\todo{embedding gemma is encoder?}
\todo{t5gemma encoder-decoder https://developers.googleblog.com/en/t5gemma/}

\todo{due parole su encoder (bert) and decoder-only}
\todo{aggiungi qualcosa come "for more/further details on xxx refer to the book jm3"}
\todo{icl, rl on human feedback, reasoning models}


\todo{transformers: encoder decoder attention self-attention. in-context learning. pretraining and finetuning transfer learning. rl tuning on human feedback. reasoning models}

\todo{large language models}



\todo{formalization, transformers, llms, bert}


\todo{Spacy architectures?. no just introduce trasformers abstractly. then spacy and others in IE}

\subsection{Large Language Models and Knowledge Bases}
\todo{blend together with llms}
\todo{comparison llms and kbs. look for the paper. lms as kbs. one may think that kbs are useless now...}
\todo{also speak of the big amount of knowledge of llms and of the research to update that knowledge. move at the beginning. deepseek r1 size. }
Modern large language models (\acsp{llm}), such as GPT~\cite{gpt4_technical_report_2023,gpt_lm_few_shot_learners_2020} and Llama~\cite{grattafiori2024llama3herdmodels}, have demonstrated remarkable capabilities in understanding natural language, enabling the application of \ac{nlp} tasks without the need for task-specific finetuning~\cite{gpt_lm_few_shot_learners_2020,instruct_gpt_2022,grattafiori2024llama3herdmodels}.
There are attempts to use them for complex tasks, including scientific discovery~\cite{ai4science2023impactlargelanguagemodels,nathani2025mlgymnewframeworkbenchmark} and intelligence analysis~\cite{fabula_intelligence_llm_2024}\todo{maybe be good to add some legal attempts to replace judges or lawyers}.
Futhermore, \acl{llm} retain considerable factual knowledge from their pretraining corpora~\cite{chang_how_do_llm_acquire_factual_knowledge_2024,physics_of_lm_3.1_2024,petroni-etal-2019-language-models-kbs,roberts-etal-2020-how-much-kg-in-llm}, and state-of-the-art models are trained on trillions of tokens; DeepSeek-V3 counts 671 billion parameters and is trained on 14.8 trillion tokens~\cite{deepseekai2025deepseekv3technicalreport}.

One may wonder why \acp{kb} are still relevant in the era of \acp{llm}. Different studies have focused the knowledge of \aclp{llm} and how they compare to \aclp{kb}~\cite{petroni-etal-2019-language-models-kbs,heinzerling-inui-2021-language}, with the survey \textcite{kg_llm_roadmap_2024} defining a roadmap for their sinergy.

% llm cons editing
While \acp{llm} possess broad knowledge, they also suffer from the incompleteness problem, as their number of parameters is limited and information is constantly evolving. Furthermore, their parametric knowledge (i.e., the knowledge stored in the model parameters) is difficult to update. There are techniques to update models' knowledge without re-train~\cite{yao-etal-2023-editing,kg_editing_survey_2024}, but with the risk of side effects, such as knowledge inconsistencies~\cite{li2024unveiling,cohen-etal-2024-rippleeffects}, with \textcite{cohen-etal-2024-rippleeffects} obtaining more consistent results with an in-context editing baseline, that adds the edited information in prompt.
% llm pro easy to learn from unstructured without humans
However, whilst \ac{kb} updates require structured facts to be inserted, often by humans~\cite{}\todo{verify humans are required}, \acp{llm} can extract knownledge from unstructured data during the training process, benefiting from data diversity (e.g., paraphrasing)~\cite{physics_of_lm_3.1_2024}.
% llm cons hallucinations
Nevertheless, \acp{llm} can generate unfaithful information, known as \emph{hallucination}~\cite{ji_etal-2023-survey-hallucination}.
% kg pro answering global questions like counts, intersections
\acp{kb} also are more interpretable, in the sense that it is possible to trace back the result of a query to the \ac{kb} facts and they do not fail global questions (e.g., ``How many presidents of the US were born in Hawaii?'') that are difficult for \acp{llm}~\cite{physics_of_lm_3.1_2024}.
\ac{kb} are consistent while \acp{llm} are not consistent and reliable~\cite{zheng2024reliablellmsknowledgebases}.\todo{verify paper}

Advantages of LLMs:
\begin{itemize}
    \item Cons: do llm know what they know? hallucinations
    \item Cons: updates (see roadmap arxiv)
\end{itemize}
Advantages of Kbs:
\begin{itemize}
    \item
\end{itemize}
\todo{little review of the combination between kbs and llms. also rag as llm combined with unstructured data (not kr or kb but still knowledge like the web)}


or \todo{kg roadmap. works that try to inject kb knowledge into llms,... or work that allow llms to query kbs? kb-rag? need also to detail rag}
\todo{final thought: kbs and llms are complementary and it is worth to work for their sinergy}
\todo{show that some paper add structure to document collections: see graphrah and co}
\todo{allucinations. llm really know what they know? find paper}

% in terms of knowledge
\Aclp{llm} posses broad knowledge as they scale is already over millions of parameters and are trained on massive corpora~\cite{}. But similarly to \acp{kb} they are not omniscient, since information is constantly evolving\todo{see above}.


Techniques to update inner \acp{llm} knowledge exist~\cite{}\todo{rome,...} but...




%
However, processing large volumes of text with \acp{llm} remains challenging\cite{}.
% long context and linear attention
Recent advances have extended the context window of \acp{llm} to millions of tokens~\cite{geminiteam2024gemini15unlockingmultimodal,kuratov2024searchneedles11mhaystack}, while other approaches aim to reduce the quadratic computational cost of full attention to linear-time complexity with respect to the number of input tokens~\cite{bigbird_2020,beltagy2020longformerlongdocumenttransformer,ding2023longnetscalingtransformers1000000000}.
But including very large texts in the context window of an \ac{llm} implies challenges such as the ``lost in the middle'' problem~\cite{liu-etal-2024-lost}---where relevant information located far from the beginning or end of a long context may be overlooked by the model---and, even if the complexity is linear, the computational cost may remain significant.
% rag and external memory
\todo{rag external memory. see also gemini 1.5 paper}
For this reason, recent work have proposed the usage of entities and knowledge graph combined with \acp{llm} to improve their performance on answering complex questions~\cite{edge2025localglobalgraphrag,hybridrag_kg_vector_rag_2024,Xu_2024_rag_kg}.\todo{read better the papers and find others. also look at matteo talk for ailc cattolica}
For instance, \textcite{edge2025localglobalgraphrag} automatically creates a \acl{kg} from a collection of documents by extracting entities and relationships, it detects communities organizin them in hierarchical levels and builds a summary for each community, allowing an \ac{llm} to answer global questions, such as ``What are the key topics represented in the dataset?''.
\todo{other usage of entities are... see matteo talk. should be good something for the retrieval}
Furthermore, semantic search capabilities, such as faceted search, can be useful tools not just for humans but also for agentic \acp{llm}~\cite{thoppilan2022lamdalanguagemodelsdialog,parisi2022talmtoolaugmentedlanguage,schick2023toolformerlanguagemodelsteach,yao2023react,gao-etal-2025-efficient}\todo{prendi da refactx qualcosa}, which can leverage them to access relevant information more effectively.

So \acl{sta} is still useful.\todo{remove} 
% graphrag
% kg RAG
% entity-enhanced retrieval
% problem of the lost-in-the-middle (https://aclanthology.org/2024.tacl-1.9/ and In search of needles in a 11m haystack)
\todo{lehmann beyond cita qualcosa di interessante per le entità con llm per taggare documenti}
\todo{LLMs cannot contain all the knowledge. different kind of memories?. RAG,...paper citato da lehmannbeyond e altri. semantic/faceted search can be useful also to llms with tools}

\section{Artificial Intelligence}
What is it, how the term is used for and maybe a definition for this thesis

Before desribing the literature of \ac{ie} and related tasks, we first introduce the specific domains that motivate this work.
\todo{improve. now that we know about incompleteness: information extraction? no better speak about the domains and them IE with examples on specific domains}